\section{Experimental design}\label{sec:design}

\subsection{Hypotheses}

The hypotheses are stated here, before the results, with the condition that refutes each; Table~\ref{tab:hyp} lists them. H1, H5, H6 and H7 are the thresholds that the sequence was designed against: a transient ratio of $0.9$ on the acceptance task, on longer moves with the frames of Eq.~\ref{eq:frames}, and the rationale of its two components, the lock and the per-frame optimization. They were fixed with the method and before the comparison with the per-frame sequence existed. H2 to H4 compare the sequence with the per-frame Gerchberg--Saxton sequence of Kim et al.~\cite{kimoe} and were stated after the sequence was designed, when that method was added to the comparison; they are post hoc with respect to the design of the method, and their thresholds come from the benchmark rule of Section~\ref{sec:stats}. The comparison with the linear time profile was added after the results and is not one of the hypotheses: it is reported as a finding made after the results (Table~\ref{tab:profile}). H8 is the consequence of the frame-count equation for a move given fewer frames than the bound asks.

\begin{table}[t]
\centering\small
\caption{Hypotheses with the condition that refutes each. The comparison hypotheses H2 to H4 are against per-frame Gerchberg--Saxton with phase induction at 5 iterations per frame; the outcome of each is in Table~\ref{tab:outcome}.}
\label{tab:hyp}
\begin{tabular}{@{}l p{6.7cm} p{6.0cm}@{}}
\toprule
ID & Statement & Refuted if\\
\midrule
H1 & On the acceptance task the sequence keeps the lowest moving trap at $0.9$ of its mean or more and above linear phase interpolation & a seed below $0.9$, or the paired interval of the difference includes zero\\
H2 & The largest phase step between frames of the sequence is smaller than that of per-frame Gerchberg--Saxton with phase induction & any seed in which the sequence has the larger step\\
H3 & The transient minimum ratio of the sequence is not lower than that of per-frame Gerchberg--Saxton with phase induction & the benchmark rule reads the dip of the published method as better\\
H4 & The largest per-frame spread of the sequence is not larger than that of the per-frame method & the benchmark rule reads the spread of the published method as better\\
H5 & With the frames of the bound, the sequence keeps the lowest moving trap at $0.9$ or more on a majority of seeds of the long moves and of the sort & a majority of seeds below $0.9$ on a move\\
H6 & Removing the trajectory lock brings the lowest moving trap below $0.9$ & the variant stays at $0.9$ or more on a majority of seeds\\
H7 & Removing the per-frame optimization brings the lowest moving trap below $0.9$ & the variant stays at $0.9$ or more on a majority of seeds\\
H8 & With fewer frames than the default, the transient ratio of the sequence falls below $0.9$ & the median at 7 frames is $0.9$ or more\\
\bottomrule
\end{tabular}
\end{table}

\subsection{Data and tasks}

All data are generated by the engine from seeds; no measured hologram or camera image enters the paper. The setup is the one of Section~\ref{sec:problem}. The acceptance task is a $4\times4$ lattice of equal traps eight bins apart in which the two right-hand columns, eight traps, move four bins in $x$; it is run at \parFrames{} frames, the default, and at $7$ and $25$ frames, with \nSeedsTwo{} seeds each. The long lattice moves are the same lattice with the two columns moved \longABins{} and \longBBins{} bins on a 64-pixel SLM, run at the frames of the bound (\longAFrames{} and \longBFrames{}) with \longSeeds{} seeds, and, where a seed fails, at twice that count. The sort is the rearrangement case of the studio: a $6\times6$ array in which \cTenMoving{} of \cTenTraps{} targets are reached by a move, with the assignment of the Hungarian method of Lee et al.~\cite{lee2017}; its longest move is \cTenFu{} Fourier units and it is run at \cTenEqual{} frames and at the \cTenFrames{} frames of the bound with \cTenSeeds{} seeds. The assignment returned by the rearrangement method differs from the plan of the case on ties, so every sequence of the sort runs on the assignment that the rearrangement method returns, and all methods move the same traps along the same paths. The ablation uses the acceptance task at \parFrames{} frames and \ablTN{} seeds.

\subsection{Methods compared}

Three sequences are run on every task. The sequence under test is the flicker-bounded sequence of Section~\ref{sec:flicker-method} with \parFrameIters{} iterations per frame. The first baseline is linear phase interpolation as described by Knottnerus et al.~\cite{lpi}: two phase-fixed weighted Gerchberg--Saxton holograms of \setupCold{} iterations for the endpoints, the position, amplitude and optical phase of each trap read from their fields and interpolated linearly over the frames (the phase by its wrapped difference), and the transform inverted at every frame without optimization. It is our implementation of a published description, not the authors' code. The second baseline is the per-frame sequence of Kim et al.~\cite{kimoe}, implemented in the engine as an option of the rearrangement method and shown as Algorithm~\ref{alg:kim}: frame 0 is the cold run of \setupCold{} iterations, and every later frame starts from the phase of the previous frame and runs a fixed number of weighted iterations toward the trap positions of that frame. It is run at \parKimIters{} iterations per frame, the count of the numerical experiment of Kim et al., and at \parFrameIters{}, the count of the proposed sequence, so that the two are also compared at equal work per frame. The weights of the feedback restart at 1 in every frame in this implementation; only the phase is carried. The paper's phase-change reduction is not implemented, because the authors did not use it in their experiments either. Both baselines use the linear constant-velocity trajectory of the rearrangement method. The proposed sequence is run with two time profiles, stated as a parameter of every comparison: its default, the minimal-jerk profile, and the linear profile of the baselines; Table~\ref{tab:profile} reads both against every baseline by the rule below. Two deviations from the paper of Kim et al.\ remain in the baseline. The zero-padding factor of the setup is \setupPad{}, while their Fig.~3 uses 4. Their frame 0 is a weighted Gerchberg--Saxton run from a random-superposition start, with a random stream of its own, whereas frame 0 of the proposed sequence is the phase-fixed rule~\cite{kimol}.

\begin{algorithm}[t]
\caption{Per-frame Gerchberg--Saxton with phase induction (Kim et al.), as run by the engine, for frames at the trap sets $X(0),\dots,X(S)$ on a linear trajectory.}
\label{alg:kim}
\begin{algorithmic}[1]
\Require traps $X(s)$ for $s=0,\dots,S$, iterations $K_0$ for frame 0, iterations $K$ per later frame, seed
\Ensure holograms $\phi^{(0)},\dots,\phi^{(S)}$
\State $\phi^{(0)}\gets$ weighted GS of $K_0$ iterations on $X(0)$ from a random-superposition start
\For{$s=1,\dots,S$}
  \State $\phi\gets\phi^{(s-1)}$ \Comment{induction: the previous phase is the initial phase}
  \State $w_m\gets1$
  \For{$t=1,\dots,K$}
    \State $V_m\gets V_m(\phi)$ at $x_m(s)$; $w_m\gets w_m\langle r\rangle/r_m$; $\phi\gets\arg\sum_mw_ma_m\,e^{i\arg V_m}\,e^{+i\Delta^m}$
  \EndFor
  \State $\phi^{(s)}\gets\phi\bmod2\pi$
\EndFor
\State \Return $\phi^{(0)},\dots,\phi^{(S)}$
\end{algorithmic}
\end{algorithm}

The difference between Algorithm~\ref{alg:kim} and Algorithm~\ref{alg:flicker} is in two lines: the phase of each coefficient is the current field's phase, free to move, and the weights restart, against the prescribed phase and the carried weights of the proposed sequence. All three sequences are scored by one function that computes Eqs.~\ref{eq:transient} to~\ref{eq:sigmaframe} from the frames and the trajectory, so that no sequence is scored on its own model.

\subsection{Metrics and statistics}\label{sec:stats}

The metrics are the transient minimum ratio $\rho$ of Eq.~\ref{eq:transient}, the per-frame minimum ratio, the largest phase step $\Theta$ of Eq.~\ref{eq:phasestep}, the largest per-frame spread $\sigma_{\max}$ of Eq.~\ref{eq:sigmaframe}, the number of frames against the frame-count equation, and the seconds of a run. The transient ratio of a method is summarized by the mean, the sample standard deviation and the interval of the mean over $n$ seeds,
\begin{equation}
\bar\rho\pm t_{n-1}\frac{s}{\sqrt n},\qquad s^2=\frac1{n-1}\sum_i\big(\rho_i-\bar\rho\big)^2,\label{eq:ci}
\end{equation}
with $t_{n-1}$ the two-sided 95 percent quantile of Student's $t$ distribution, and the comparison of the sequence with a baseline on the same seed index is the paired difference
\begin{equation}
\Delta_i=\rho_i^{\mathrm{seq}}-\rho_i^{\mathrm{base}},\qquad \bar\Delta\pm t_{n-1}\frac{s_\Delta}{\sqrt n}.\label{eq:paired}
\end{equation}
The pairing is by seed index only: the proposed sequence draws its random start from the seed, and the baselines draw theirs from an offset seed, so the methods do not share a random start and the paired interval is not a matched-start comparison. With $n$ of six or eight the intervals are wide and the tests that they imply have low power; no $p$-value is reported, and a claim rests on the interval, the count of seeds that meet a threshold and the reading rule below. Medians and ranges are reported as well because the distributions are not normal.

The reading of a comparison, which method is better on a measure, is the rule of the studio's benchmark and is applied identically to every cell: at least \winMinSeeds{} seeds on both sides, interquartile ranges that do not overlap, and a gap of at least \winMinDiff{} percentage points of the dip below the mean, $100(1-\rho)$, which is lower for the better method, or of the per-frame spread; otherwise the cells are read as a tie. The word \emph{proposed} in a table reads the sequence as better by this rule, \emph{published} reads the baseline as better, and \emph{tie} reads neither. The rule is conservative: a difference in the mean can be real and read as a tie.

\subsection{Protocol}

Every method is deterministic given a seed through the portable generator of the engine, so a rerun reproduces the numbers up to the floating point of the machine's FFT; the seconds differ between runs. The seeds are $1,\dots,n$ with $n$ stated per experiment in the previous subsections. The seconds are wall-clock times of one run of the sequence on the machine that produced the data, with another job loading the machine during part of the runs, so they are reported as an order of magnitude and not as a benchmark. The ablation removes one component at a time and keeps the others, on the same seeds and task as the main result, and the variant with all components is checked against the module of the engine (Section~\ref{sec:flicker-method}). The tables of the results are written by the same script that runs the experiments, from the stored measurement file, so that a table cannot differ from the figure drawn from the same file. Every comparison reports the number of seeds on each side, since the reading rule needs at least \winMinSeeds{} on both.
